% !TEX TS-program = xelatex
% !TEX encoding = UTF-8 Unicode
% ================================================
%  上海交通大学 研究生工作周报
%  编译: xelatex → bibtex → xelatex → xelatex
% ================================================

\documentclass{sjtuweekly}
\usepackage{booktabs}
\usepackage{tabularx}
\usepackage{array}
\usepackage{enumitem}

\renewcommand{\thefigure}{\arabic{figure}}
\renewcommand{\thetable}{\arabic{table}}

\bibfile{bib/database}

\member{霍锦龙}
\weeknum{xx}

\begin{document}

\makeweeklytitle

% ================================================
\section{一、本周工作概要}

\subsection*{1.1\quad 实验：LoRA 微调 MoE Router Gate}

在 Qwen3.6-35B-A3B: 256 专家 × 40 层，top-8 上对 router gate 施加 LoRA rank=16，20 条 MoE 合成数据，对比训练前后同一输入的路由分布：

\begin{equation}
  W_{\text{gate}} \leftarrow W_{\text{gate}} + \alpha \cdot B \cdot A,\quad
  S(x) = \text{top-}k\big(\text{softmax}(x \cdot W_{\text{gate}})\big)
  \label{eq:lora_gate}
\end{equation}

\begin{itemize}
  \item MLX / HF 双后端路由捕获输出 \texttt{RoutingTrace} JSON
  \item 分析训练中的负载均衡/熵/共激活、\texttt{compare}、\texttt{intervene}、\texttt{ablate}
\end{itemize}


\subsection*{1.2\quad OCS 预触发逻辑设计}
如之前的讨论和RP中的内容,如果为ocs提供moe的专家偏好,首先要捕获模式,并进行预配置和CCL相关匹配.

\textbf{Pipeline思路/逻辑：} LoRA 训练 → 捕获路由 → 构建协同激活图 → OCS 预配置表 推理时 GPU 间电路预连线，命中率 $\sim$71\%，miss 回退动态切换。

\textbf{数据支撑：} 活跃专家 -42.2\%，协同激活对 -43.3\% 更少的 OCS 电路需管理。


\subsection*{1.3\quad 发现}

\begin{enumerate}
  \item \textbf{路由可学习且方向稳定:} 同一输入，训练前后 71.1\% 路由决策一致:层 0: 93.4\% → 层 39: 54.5\%）。浅层保留、深层变换，符合"语法→语义"分层假说。且可以被用于推理,Pagedattention的设计思路.
  \item \textbf{路由集中化:} 活跃专家数大幅下降, 从1到层38: 专家激活从161到68，无死专家256/256 均被某层使用）。专家 41 激活次数最多有 32 次分配。
  \item \textbf{协同激活:} Top-20 共激活对训练后全部持续存在，41, 151, 224 始终作为适配的激活专家。

\end{enumerate}


% ================================================
\section{二、遇到的问题与讨论}

\subsection*{2.1\quad 统计方法问题}

\begin{enumerate}[label=(\arabic*), leftmargin=*]
  \item \textbf{构造错配:} 当前的"71\% 训练前后一致"基于比较脆弱的假设,即当前的数据样本比较少.同时考虑到有相关文献指出应用于 OCS 需要算法或验证模型具备一定预测性/泛化性.即 gate 计算前预测未见 token 的专家集。从这个角度出发, 严格来说,两者数学上不同，前者无法推导后者。但是,在当前的预训练/后训练下,刨除多模态moe的影响,单一模态具备类似数学上的预测性.也就是具备一定程度上的可推广性.
  \item \textbf{n=1 无方差:} 全部数据来自单条 36-token 序列。偏好专家可能是该 prompt 下的偶然偏好，无法外推。
  \item \textbf{训练退化:} checkpoint 有时会生成崩溃，没有涉及"退化崩溃前后"专家下降的贡献度讨论。
\end{enumerate}

\subsection*{2.2\quad 系统建模问题}

\begin{enumerate}[label=(\arabic*), leftmargin=*]
  \item \textbf{硬件粒度错配:} 活跃专家总数下降不减少单 token 电路需求。例如 Top-8 每 token 始终需 8 专家. 那么后续就要考虑在OCS场景下的并发token在同一时间窗内的专家集重叠度和资源收集以及调度问题。
  \item \textbf{缺乏OCS延迟预算:} MEMS OCS 切换 $\sim$25ms，单层推理微秒级, $\tau_{\text{reconf}} \gg \tau_{\text{compute}}$ 时单 token 粒度预触发物理不可行。需明确硬件假设，或粗化到批次粒度（MixNet 思路）,这周的工作内容里还未和ocs进行.
  \item \textbf{MoE流量矩阵和OCS重配置适配度:} 分布式 MoE 是 All-to-All 集合通信. 部分文章中指出 OCS 在一定程度上要支持批次内的置换矩阵，而非单个专家对之间的连线(性价比不高)。后续可能要考虑 A2A 流量矩阵的稀疏度来决定批次内训练/推理的拓扑设置。

\end{enumerate}

% ================================================
\section{三、工作计划}

\begin{table}[htbp]
\centering
\small
\begin{tabularx}{\linewidth}{@{}XX@{}}
\toprule
内容 & 预期产出 \\
\midrule
\textbf{多样化微调:} 使用跨领域样本再次微调，确保生成数据分布和路由偏好不退化，复测活跃专家/熵是否仍下降 & checkpoint 失效前后的分布讨论 \\
\textbf{硬件建模:} 选定 $\tau_{\text{reconf}}$（MEMS/photonic/SOA）和 $\tau_{\text{compute}}$（目标 batch），计算期望延迟公式，判定物理可行性 & 项目是否有物理意义 \\
\textbf{构建预测器:} 用前一层的 expert set 或 pre-gate activations 做线性 probe 预测当前层 expert set，对比 chance baseline（3.1\%）和 SOTA（93--97\%） & 验证MoE专家网络token可预测性的上下界 \\
\textbf{系统对标:} 切换分析单元到并发 token 专家集重叠度；与 MixNet / pre-attention predictor 并列对比；多租户混合流量压力测试 & 在系统社区中的定位 \\
\bottomrule
\end{tabularx}

\caption{后续计划}
\label{tab:plan}
\end{table}


\printbib

\end{document}
